\chapter{Lecture 16: Abstract Classes, Templates, Vectors and Exceptions}

\tikzset{umlc/.style={draw,fill=blue!4,rectangle split,rectangle split parts=2,
    rectangle split part align={center,left},font=\small,inner sep=3pt,minimum width=2.6cm},
  inh/.style={-{Triangle[open,length=3mm,width=3mm]}}}
\begin{center}
\textit{March 6, 2018}
\end{center}

\section{Pure Virtual Methods and Abstract Classes}
\begin{lstlisting}[style=cpp]
class Student {
    ...
  public:
    virtual int fees() const = 0;    // pure virtual method
    virtual ~Student() {}
};

class Regular : public Student {
  public:
    int fees() const override { ... }  // regular fees
};

class CoOp : public Student {
  public:
    int fees() const override { ... }  // co-op fees
};
\end{lstlisting}
What should \code{Student::fees} do? There's no sensible answer: every student
is either regular or co-op. So we don't implement it.

\begin{definition}
A \textbf{pure virtual method} (declared with \code{= 0}) is a method that need
not have an implementation in its class. A class with at least one pure virtual
method is an \textbf{abstract class}, and cannot be instantiated. A class that
is not abstract is \textbf{concrete}.
\end{definition}
\begin{lstlisting}[style=cpp]
Student s;                    // won't compile: Student is abstract
Student *sp = new CoOp{...};  // fine
\end{lstlisting}
A subclass inherits all members, including pure virtual methods; so a subclass
must implement \emph{all} inherited pure virtual methods to be concrete,
otherwise it is abstract too.

Uses of abstract classes:
\begin{itemize}
  \item organize subclasses;
  \item hold the fields and methods common to all subclasses;
  \item serve as the type of pointers/references, so polymorphism works: a
        \code{Student*} points to either a Regular or a CoOp student (never to
        a plain Student, since none can exist).
\end{itemize}

\subsection*{UML}
In UML, virtual and pure virtual methods are written in \emph{italics}, and so
are the names of abstract classes. Static members are \underline{underlined}.
(Tools: StarUML, OmniGraffle, \dots)

\begin{center}
\begin{tikzpicture}
  \node[umlc] (s) at (0,0) {\textit{Student}\nodepart{second}\code{+}\textit{fees(): Integer}};
  \node[umlc] (r) at (-2.2,-2.2) {Regular\nodepart{second}\code{+fees(): Integer}};
  \node[umlc] (c) at (2.2,-2.2) {CoOp\nodepart{second}\code{+fees(): Integer}};
  \draw[inh] (r.north) -- (s.south west);
  \draw[inh] (c.north) -- (s.south east);
\end{tikzpicture}
\end{center}

\section{Templates}
\begin{lstlisting}[style=cpp]
class Stack {
    int size, cap;
    int *contents;
  public:
    Stack();
    void push(int x);
    void pop();
    int top() const;
    ~Stack();
};
\end{lstlisting}
This works only for \code{int}s. A stack of \code{string}s or of
\code{Expression*}s would be the same code with a different type.

\begin{definition}
A \textbf{template class} is a class parameterized on one or more types.
\end{definition}
\begin{lstlisting}[style=cpp]
template <typename T> class Stack {
    int size, cap;
    T *contents;
  public:
    Stack();
    void push(T x);
    void pop();
    T top() const;
    ~Stack();
};

Stack<int> sInts;           // template instantiation with T = int
sInts.push(1);
int x = sInts.top();

Stack<Expression*> sExpr;
sExpr.push(new BinExp{...});
\end{lstlisting}
The compiler takes the template and generates (stamps out) a separate copy of
the class for each type it is used with, substituting the type for \code{T}.
Because of this, the full template implementation must be visible where it is
used: templates are written entirely in the header file.
\begin{center}
\begin{tikzpicture}[>=Stealth,font=\small]
  \node[draw,dashed,rounded corners,fill=gray!8,minimum height=0.9cm,align=center] (t) at (0,0) {\code{template <typename T>}\\ \code{class Stack}};
  \node[draw,fill=blue!5,minimum height=0.8cm] (a) at (-2.6,-1.8) {\code{Stack<int>}};
  \node[draw,fill=blue!5,minimum height=0.8cm] (b) at (2.6,-1.8) {\code{Stack<Expression*>}};
  \draw[->] (t) -- node[left,font=\scriptsize]{\code{T = int}} (a);
  \draw[->] (t) -- node[right,font=\scriptsize]{\code{T = Expression*}} (b);
\end{tikzpicture}

\small One class is generated per type the template is used with.
\end{center}

\subsection{A Template List with an Iterator}
\begin{lstlisting}[style=cpp]
template <typename T> class List {
    struct Node {
        T data;
        Node *next;
        Node(const T &data, Node *next) : data{data}, next{next} {}
        ~Node() { delete next; }
    };
    Node *theList = nullptr;
  public:
    class Iterator {
        Node *curr;
        explicit Iterator(Node *curr) : curr{curr} {}
      public:
        T &operator*() const { return curr->data; }
        Iterator &operator++() { curr = curr->next; return *this; }
        bool operator!=(const Iterator &other) const {
            return curr != other.curr;
        }
        friend class List<T>;
    };
    Iterator begin() { return Iterator{theList}; }
    Iterator end() { return Iterator{nullptr}; }
    void addToFront(const T &n) { theList = new Node{n, theList}; }
    T &ith(int i);
    ~List() { delete theList; }
};
\end{lstlisting}
\begin{lstlisting}[style=cpp]
List<int> l1;
l1.addToFront(1);

List<List<int>> l2;     // a list of lists of ints
l2.addToFront(l1);      // (needs a correct copy ctor for List!)

for (auto &m : l2) {
    for (auto n : m) cout << n << endl;
}
\end{lstlisting}

\section{The STL: Dynamic-Length Arrays}
The \textbf{Standard Template Library} provides many template classes. The most
useful is \code{std::vector}:
\begin{itemize}
  \item guaranteed to be implemented as an array (elements are contiguous);
  \item automatically resizes as needed.
\end{itemize}
\begin{lstlisting}[style=cpp]
#include <vector>
using namespace std;

int main() {
    vector<int> v{4, 5};   // [4, 5]; v is on the stack, its array on the heap
    v.emplace_back(6);     // [4, 5, 6]
    v.push_back(7);        // [4, 5, 6, 7]
    vector<int> w(4, 5);   // round brackets: 4 elements, each 5: [5, 5, 5, 5]
}
\end{lstlisting}
Note the difference between \code{\{4, 5\}} (the elements) and \code{(4, 5)}
(count and value).

\begin{center}
\begin{tikzpicture}[>=Stealth,font=\small,x=0.8cm,c/.style={draw,minimum width=0.8cm,minimum height=0.65cm}]
  \node[draw,fill=blue!5,minimum width=2.4cm,minimum height=0.9cm,align=center,font=\scriptsize] (v) at (-3,0) {\code{v}: data pointer,\\ size 4, capacity};
  \node[font=\scriptsize] at (-3,0.75) {stack};
  \foreach \x [count=\k from 0] in {4,5,6,7} \node[c,fill=green!8] (e\k) at (\k,0) {\x};
  \node[c,draw=gray] (e4) at (4,0) {};
  \node[c,draw=gray] (e5) at (5,0) {};
  \node[font=\scriptsize] at (2.5,0.75) {heap array (spare capacity at the end)};
  \draw[->] (v.east) -- (e0.west);
  \draw[<-] (e0.south) -- ++(0,-0.6) node[below,font=\scriptsize]{\code{v.begin()}};
  \draw[<-] (e4.south) -- ++(0,-0.6) node[below,font=\scriptsize]{\code{v.end()}};
\end{tikzpicture}

\small \code{end()} points one past the last element, so the loop condition \code{it != v.end()} stops on time.
\end{center}

Ways to loop over a vector:
\begin{lstlisting}[style=cpp]
for (size_t i = 0; i < v.size(); ++i) cout << v[i] << endl;  // operator[]

for (vector<int>::iterator it = v.begin(); it != v.end(); ++it)
    cout << *it << endl;

for (auto n : v) cout << n << endl;

for (vector<int>::reverse_iterator it = v.rbegin(); it != v.rend(); ++it)
    cout << *it << endl;                     // backwards
for (auto it = v.rbegin(); it != v.rend(); ++it) ...
\end{lstlisting}
Highly recommended: read the reference for \code{vector}.

\begin{lstlisting}[style=cpp]
v.pop_back();             // remove the last element
v.erase(v.begin());       // remove the first element: O(n)
v.erase(v.end() - 1);     // remove the last element
\end{lstlisting}
The STL makes heavy use of iterators. \code{vector::erase} takes an iterator,
and returns an iterator to the element after the one just erased (iterators to
erased or later elements are invalidated).

\begin{important}[Assignment 4]
A4 prohibits heap arrays (\code{new Abc[n]}): use \code{vector} instead. A vector
manages its own memory, so classes using it need no Big Five.
\end{important}

\subsection{Checked and Unchecked Access}
\code{v[i]} accesses the $i$th element \emph{unchecked}: the program does not
check whether \code{i} is in bounds, and an out-of-bounds access is undefined
behaviour. \code{v.at(i)} is \emph{checked} access (Java does checked access by
default).

\paragraph{Q.} What happens if the argument to \code{at} is out of range?

\paragraph{A.} An exception is thrown (raised).

\section{C++ Exceptions}
\begin{lstlisting}[style=cpp]
// client code
vector<int> v;
v.at(5);     // vector::at detects the error -- but it is the client
             // who knows how to deal with it
\end{lstlisting}

\paragraph{Q.} How should \code{vector::at} report the error to the client?

In C:
\begin{itemize}
  \item return an error value (e.g.\ $-1$) --- but every \code{int} might be a
        legitimate element;
  \item set a global error variable, e.g.\ \code{errno}.
\end{itemize}
Both are bad:
\begin{itemize}
  \item awkward to code with \code{if} statements after every call;
  \item programmers are lazy, so errors go undetected.
\end{itemize}
\begin{quote}
\textit{``You are to stop being lazy right off the bat.''}
\end{quote}

In C++, when an error condition occurs, the function \textbf{throws an
exception}. Control immediately leaves the function and goes back up the call
chain looking for a \emph{handler}. If nobody handles it, the program
terminates. So the error can't be silently ignored.

\begin{lstlisting}[style=cpp]
#include <stdexcept>
vector<int> v;
try {
    cout << v.at(10000) << endl;   // throws std::out_of_range
    cout << "never printed" << endl;
} catch (out_of_range &r) {
    cerr << "Range error: " << r.what() << endl;
}
\end{lstlisting}
\code{what()} returns a string describing the error. The details follow next
lecture.

\section*{Summary}
\begin{itemize}
  \item Pure virtual (\code{= 0}) methods make a class abstract; subclasses must
        implement them all to be concrete.
  \item Templates parameterize classes on types; they live in headers.
  \item \code{std::vector}: a self-resizing array; iterate with indices,
        iterators, range-for or reverse iterators; \code{erase} takes iterators.
  \item \code{v[i]} is unchecked; \code{v.at(i)} throws \code{out\_of\_range}.
  \item Exceptions replace error codes: they can't be ignored, and are caught
        with \code{try}/\code{catch}.
\end{itemize}
